\subsection{Voltage Regulator (LM7805CV)}
\markright{Voltage Regulator}

\subsubsection*{Definition}
A linear voltage regulator is a three-terminal IC that maintains a fixed, stable DC
output voltage regardless of variations in the input supply or load current
\cite{sedra2020}. The LM7805CV is a member of the 78xx series, providing a fixed
$+5$\,V output. Internally it uses a Zener-diode reference, an error amplifier, and a
pass transistor connected as an emitter follower to regulate the output.

\labimage{lab-01/assets/lm7805.jpeg}%
  {LM7805CV $+5$\,V linear voltage regulator in TO-220 package. Pins
   from left: input, ground, output.}{lm7805_lab}

\subsubsection*{Specifications}
The LM7805 accepts input voltages from 7\,V to 35\,V and delivers a regulated $+5$\,V
output at up to 1\,A continuous current ($1.5$\,A with heatsink). Line regulation is
typically 3\,mV (output change for full input swing); load regulation is about 15\,mV
\cite{boylestad2015}. A 0.33\,$\mu$F input decoupling capacitor and 0.1\,$\mu$F output
capacitor are recommended for stability.

\subsubsection*{Purpose of Use}
The LM7805 converts an unregulated DC supply (e.g., a 9\,V battery or 12\,V rectified
supply) into a stable $+5$\,V rail for powering digital logic circuits, microcontrollers,
and sensor modules.

\subsubsection*{Applications}
The 78xx family is used in laboratory bench power supply circuits, embedded system
prototypes, battery-powered sensor nodes, and educational PCB designs \cite{floyd2012}.
